\chapter*{Abstract}

We study adaptive target detection in unknown and possibly non-Gaussian background clutter using target-free secondary data. We take a hybrid model-based and data-driven approach. Learning is used to model the background, while the detector and its normalization are kept in the form prescribed by classical detection theory. Our focus is on the locally most powerful (LMP) detector that depends on the background only through its score function and a few score moments. We exploit this structure by estimating the background characteristics with denoising score matching (DSM) and plugging the estimate into the detector. This separation allows a single learned background model to be reused across target signatures. In the affine Gaussian case, the proposed method reduces to a regularized adaptive matched filter, which explains its advantages in finite sample regimes. We also develop a spatially adaptive extension that improves accuracy while maintaining empirical constant false alarm rate (CFAR) calibration in spatially varying environments.
Experiments on three real hyperspectral scenes illustrate the advantages of the proposed detectors with respect to Gaussian, Gaussian mixture, learned Rao, and recent deep learning baselines. Code and complete results: \url{https://github.com/michaelpiro/DART}.